\section{Beneficiaries}

The research results and the technical framework developed in this study are intended to benefit multiple stakeholders in Vietnam's artificial intelligence and legal sectors:

\begin{itemize}[leftmargin=*]
    \item \textbf{Legal NLP Research Community:} The proposed benchmark and quantitative metrics will establish a standardized baseline for evaluation of continual learning in low-resource Vietnamese legal systems. Researchers can utilize the open-source code and metadata structures to test alternative retrieval models and memory paradigms.
    \item \textbf{Legal Technology Developers:} Software developers constructing legal assistant systems, court case search platforms, or statutory QA tools can directly reuse the selective memory and unlearning gates. This reduces the development overhead of ensuring legislative compliance and security governance.
    \item \textbf{Regulatory and Compliance Institutions:} Organizations that require data security, privacy controls, and legislative auditing (such as banking compliance departments, Supreme Court auditing teams, and private law firms) can deploy the proposed behavioral unlearning modules. This guarantees that confidential or repealed information is never leaked or cited in legal reasoning, thereby mitigating compliance risks.
\end{itemize}
